\section{Credit schedules after financial write-off}
\label{sim:credit-claim-continuity}

Financial closure leaves a separate question about the history of scheduled
receipts. The retained credit examples already record uncollected amounts and
portions removed by write-off. A dated reconciliation now checks how those
amounts relate to the original schedule, actual collections and any remaining
modeled receipts. These quantities are nominal currency amounts. They are
neither discounted carrying values nor an assessment of enforceable legal
rights.

Let $d$ denote a scheduled payment date and $k$ the position in the recorded
event sequence. Write $O_d$ for the original scheduled amount, $\Delta_{k,d}$
for cumulative signed changes to the modeled receipt schedule, $C_{k,d}$ for
collections, $U_{k,d}$ for uncollected amounts recorded when a due date is
resolved, $W_{k,d}$ for portions removed by financial write-off, and $R_{k,d}$
for the remaining schedule. Each dated record satisfies
\[
 O_d+\Delta_{k,d}=C_{k,d}+U_{k,d}+W_{k,d}+R_{k,d}.
\]
A collection transfers the scheduled amount at that date from $R$ to $C$ and
$U$. A proportional write-off transfers the specified portion from $R$ to
$W$. A replacement schedule changes $R$ and $\Delta$ by the same dated amount.
Accrual changes the financial measurements but leaves this nominal relation
unchanged. These observations derive the identity from the admitted events;
they do not establish which events should be admitted for a real loan.

In the partial-write-off example, the original schedule is 66 on 1 January
2026 and $363/5$ on 1 January 2027. The identified one-fifth portion removed
on 1 July 2025 contains $66/5$ and $363/25$. Subsequent collections of 44 and
55 leave uncollected amounts $44/5$ and $77/25$. At the two dates separately,
\[
 66=44+\frac{44}{5}+\frac{66}{5},\qquad
 \frac{363}{5}=55+\frac{77}{25}+\frac{363}{25}.
\]
Their sum is $693/5=99+297/25+693/25$. The full-write-off example instead
retains the whole original schedule among the removed portions and keeps its
zero-valued remaining entries. The record distinguishes those entries from
missing dates.

The performing prepayment illustrates why the signed revision belongs in the
identity. Under its supplied unchanged-terms finding, the modeled schedule
changes from 55 and $121/2$ at the two annual dates to 105 at the first date.
The changes are therefore $+50$ and $-121/2$. The later receipt of 105 reconciles
to that explicit replacement. Calling the replacement a legal amendment, or
silently comparing the receipt with the old total, would change the meaning
of the example.

The read-only companion checks all 57 stages of the nine retained histories
against their original schedules and archives, with independently authored
endpoints. It retains event order even when two events share a date, zero
uncollected amounts, write-off portion identities, and the exact source and
financial-result identities. Twenty-five controls reject missing or duplicated
portions, changed dates, hidden replacements and unsupported assertions of
forgiveness or accounting admission. The existing financial journals remain
unchanged. Reproducible calculations and their limits are recorded in
\path{data/accounting-master/credit-claim-continuity-20261027/result.md}.

All nine examples finish with a financially closed book. Adding a further event
through the existing native interface produces \texttt{event\_after\_close}.
That refusal describes the implemented interface; it does not determine whether
a recovery can occur. A later recovery needs its own dated receipt evidence,
allocation and supported accounting treatment. The present record preserves
the earlier history while leaving forgiveness and later recovery unassessed.
